\documentclass[letterpaper,12pt]{article}
\setlength{\topmargin}{0in}
\setlength{\oddsidemargin}{0in}
\setlength{\evensidemargin}{0in}
\setlength{\textwidth}{6.5in}
\setlength{\headheight}{0in}
\setlength{\headsep}{0in}
\setlength{\topskip}{0.5in}
\setlength{\textheight}{9in}
\setlength{\parindent}{0.0in}
\setlength{\parskip}{1em}
\usepackage{common,epsfig}

\pagestyle{empty}

\begin{document}

\today

Grainger Search Committee\\
University of Wisconsin\\
Astronomy Department\\
475 N. Charter Street\\
Madison, WI 53706\\
USA

To whom it may concern:

Please find attached my 2011 Grainger Fellowship application titled
``Mapping Galaxy Cluster Magnetic Fields: An Observational Study of
ICM Physics.'' I feel the University of Wisconsin-Madison (UW) is an
excellent fit for my research goals and the professional environment
will benefit from my addition. The UW Astronomy and Physics
Departments, Center for Plasma Theory \& Computation, and Center for
Magnetic Self-Organization are home to numerous experts in one, or
several, of the science areas discussed in this proposal
(\eg\ X-ray/radio/polarization studies of galaxy clusters, AGN
feedback, computational modeling, high-energy astronomy, plasma
physics). Their experience with the latest techniques in all these
disciplines will be invaluable to the success of my proposed work. In
addition, my expertise in X-ray and radio astrophysics, and mature
collaborations with groups in both research areas, has prepared me for
the rigors of the program proposed in this application.

Along with this letter are my curriculum vitae, a list of
publications, and a statement of my proposed research program. Letters
of recommendation from Dr. Megan Donahue, Dr. Mark Voit, and Dr. Brian
McNamara should arrive under separate cover. Please do not hesitate to
contact me if there is any further information I can provide as you
review my application. Thank you for your consideration.

Sincerely,\\
\begin{minipage}{7.5in}
\epsfig{file=signature.ps,height=0.75in}
\end{minipage}
Kenneth W. Cavagnolo

\end{document}
